\section{Introduction}\label{sec:introduction}

A lower bound for scalar curvature is a local condition. Gromov's
macroscopic dimension conjecture predicts a global consequence: at a
controlled scale, two dimensions can be collapsed. We prove the
quantitative form of this assertion for arbitrary complete manifolds
in dimensions at least four.

For a metric space \(X\) and an integer \(d\ge0\), define its
\emph{Urysohn \(d\)-width} by
\[
 \UW_d(X)=\inf_f\ \sup_{y\in K}\diam_X f^{-1}(y),
\]
where the infimum runs over continuous maps \(f:X\to K\) to simplicial
complexes of dimension at most \(d\). The complexes have their usual
polyhedral topology. The diameter is that of the entire fiber, including
all its components; empty fibers may be omitted.

\begin{theorem}\label{thm:main}
For every integer \(n\ge4\), there is a finite constant \(C_n\) such
that every connected complete smooth Riemannian \(n\)-manifold
\((M,g)\), without boundary and with \(\Scal_g\ge1\), admits a
continuous map \(f:M\to K\) to a simplicial complex of dimension at
most \(n-2\) satisfying
\[
 \diam_g f^{-1}(y)\le C_n\qquad(y\in K).
\]

\end{theorem}

Thus \(\UW_{n-2}(M,g)\le C_n\). The constant is independent of \(M\), its topology, and its geometry.
The theorem applies to compact and noncompact manifolds. If
\(\Scal_g\ge\sigma^2>0\), applying the theorem to \(\sigma^2g\) gives
\(\UW_{n-2}(M,g)\le C_n/\sigma\).

For a closed connected Riemannian manifold, let
\(\operatorname{FillRad}(M,g)\) be Gromov's filling radius
\citep{gromov1983Filling}, defined by the Kuratowski embedding and the
fundamental class, with coefficients in
\(\Z\) if \(M\) is orientable and in \(\Z_2\) otherwise
\citep[Definition~3.1]{nabutovskyRotmanSabourau2021}.
The next corollary gives the closed-manifold case in dimensions at least
four of Gromov's filling-radius Conjecture~\(A^{-}\)
\citep[\S3]{gromov2017Dozen}.

\begin{corollary}[Closed-manifold filling-radius bound]
\label{cor:filling-radius}
Let \((M,g)\) be a closed connected smooth Riemannian \(n\)-manifold
with \(n\ge4\). If \(\Scal_g\ge\sigma^2\), where \(\sigma>0\), then
\[
 \operatorname{FillRad}(M,g)\le \frac{C_n}{2\sigma},
\]
where \(C_n\) is the constant in Theorem~\ref{thm:main}.
\end{corollary}



For the closed-cover consequence we also use the \emph{continuous
macroscopic dimension}: $\dim_{\mathrm{mc}}X\le d$ means that there is a
continuous map to a simplicial complex of dimension at most $d$, with its uniform simplex metric
and a finite uniform bound on the diameters of its entire fibers.
Each simplex has the same standard Euclidean shape and edge length.

\begin{corollary}[Continuous macroscopic dimension of universal covers]
\label{cover:main}
Let $(M,h)$ be a closed connected smooth Riemannian manifold of
dimension $n\ge2$ with positive scalar curvature.  Let $\widetilde h$ be the lift of $h$ to its universal cover
$\widetilde M$. Then
\[
 \dim_{\mathrm{mc}}(\widetilde M,\widetilde h)\le n-2.
\]
\end{corollary}

\begin{corollary}[Aspherical obstruction]\label{cor:aspherical}
A closed smooth aspherical manifold of dimension \(n\ge4\) admits
no Riemannian metric of strictly positive scalar curvature.
\end{corollary}



An independent proof in every dimension follows from the stronger
rational-inessentiality theorem of
\citet[Theorem~1.1 and Corollary~1.2]{openai2026RationalInessentiality}.

\subsection{History and significance}

Gromov introduced macroscopic dimension in his study of positive scalar
curvature \citep[\S2]{gromov1996}. His quantitative conjecture already asks
whether a complete manifold with a positive scalar-curvature lower bound
can be mapped continuously to a polyhedron of codimension two with all
point inverses bounded at the curvature scale
\citep[\S\(2\frac12\)]{gromov1996}; see also
\citep[p.~3180, Comments after Conjecture~3.2]{gromov2019Mean}.
The diameter condition concerns each entire fiber, including all its
components. In the universal cover of a closed manifold the corresponding
macroscopic-dimension assertion asks only for some finite bound.
The continuous invariant and the version requiring a Lipschitz map to a
uniformly metrized complex are distinct
\citep[Definitions~1.1--1.2, Theorem~6.3 and the following Example]{dranishnikov2014}.

The pointwise curvature hypothesis cannot be replaced by its fixed-scale
macroscopic counterpart \citep[Definition~1.3]{kumarSen2026Width}: at a
prescribed radius \(r>0\), the macroscopic scalar curvature at a point is
the scalar curvature of the simply connected space form of the same
dimension whose \(r\)-balls have the volume of the \(r\)-ball about a lift
of the point in the universal cover of the concentric \(2r\)-ball. Kumar and Sen
construct closed manifolds in
every dimension at least four with a prescribed positive fixed-scale
macroscopic scalar-curvature lower bound and arbitrarily large
Urysohn codimension-two width \citep[Theorem~A]{kumarSen2026Width}.

The low-dimensional cases have a substantial earlier history.
Gromov and Lawson obtained geometric width bounds for closed orientable
three-manifolds \citep[Corollary~10.11]{gromovLawson1983}, and Bolotov proved that the
universal cover of a closed three-manifold cannot have macroscopic
dimension two \citep{bolotov2003}.
For closed three-manifolds with positive scalar curvature, Liokumovich
and Maximo constructed Morse functions with curvature-scale bounds for
the areas and ambient diameters of connected level components
\citep[Theorem~1.1]{liokumovichMaximo2023}.
Liokumovich and Wang constructed proper Morse functions on
complete noncompact three-manifolds with uniformly positive scalar
curvature, with uniform diameter bounds for connected level components
\citep[Theorem~1.1]{liokumovichWang2023}.
We use that result and an explicit passage to the graph obtained by collapsing each connected level component to a point for the
three-dimensional closed-cover consequence; the surface case follows
from the Bonnet--Myers argument.

The aspherical obstruction in dimensions four and five was established
by Chodosh and Li \citep[Theorem~2]{chodoshLi}, with an independent
proof in dimension five by Gromov
\citep[\S7, Main Corollary]{gromov2020Aspherical}.

For the macroscopic-dimension conjecture, Bolotov proved the closed spin
case under the Strong Novikov Conjecture and a cohomological-dimension condition
\citep{bolotov2009}. Bolotov and Dranishnikov then obtained the
codimension-two conclusion for closed spin manifolds under the
Strong Novikov Conjecture and
injectivity of the connective-to-periodic real K-homology map
\citep[Theorem~4.3]{bolotovDranishnikov2010}.
For closed orientable totally non-spin manifolds of dimension at least
five, they also showed that a macroscopic dimension drop by one implies
a drop by two \citep[Theorem~3.4]{bolotovDranishnikov2016}.
Daher and Dranishnikov established the analogous implication for closed
orientable totally non-spin four-manifolds with residually finite
fundamental group
\citep[Theorem~3.3]{daherDranishnikov2019}.
These topological advances concern universal covers and finite bounds;
the present quantitative statement is formulated for arbitrary complete
manifolds with a bound depending only on dimension after scalar normalization.


The geometric method uses the minimal-hypersurface approach of Schoen
and Yau, in which stability and the Gauss equation pass scalar-curvature
information to hypersurfaces \citep[Section~1]{schoenYau1979Structure},
and the torical symmetrization mechanism described by Gromov
\citep[Section~2]{gromov2018Metric}. Gromov's \(\mu\)-bubble method
uses prescribed mean curvature to localize the hypersurfaces entering
scalar-curvature descent \citep[\S\(5\frac56\)]{gromov1996}.
The weighted variational formulation and its relation to warped-circle
geometry and stability are developed by Chodosh and Li
\citep[Section~3 and Remark~11]{chodoshLi}.

Removing singular regions while retaining scalar-curvature information
is also central to Lohkamp's splitting theory
\citep[\S1.2 and Theorem~3.15]{lohkamp2022SplittingsII}.
In the construction below, function updates exclude singular points from
the regions where the next scalar bound is required; the support of the
minimizer's perimeter measure and both retained sides are then regularized together.

Recent results also address quantitative width and singular descent.
Antonelli proves width bounds under a positive scalar-curvature lower
bound, an almost-nonnegative Ricci lower bound, and uniform noncollapse,
with constants depending on the volume lower bound
\citep[Theorem~1.4]{antonelli2026Volume}.
Bi and Zhu develop singular scalar-curvature descent for enlargeability
obstructions and cubical-width inequalities
\citep[Theorem~B and Section~4]{biZhu2026}.
Rather than continuing the descent on a single hypersurface, the present
construction keeps both sides of every cut as well as its boundary support, through
countably many cuts.

\subsection{The proof and its main constructions}

The proof of \Cref{thm:main} constructs separating sets with a bound
on their intersections. Fix a length \(D\), depending only on \(n\),
and a separated net \((x_i)\) whose closed \(D\)-balls cover \(M\)
and whose closed \(2D\)-balls are locally finite. Give each center a
label. Since the cuts insert cylinders, we record original positions
by a continuous collapse \(q:M\to M\) with \(d_g(q(x),x)<1\).
The required separating set for label \(i\) separates the \(q\)-preimages of
the inner ball and the exterior of the outer ball; no \(n-1\)
distinct separating sets may meet. These are the partitions of
\Cref{topology:partitions}. To see why they suffice, thicken the separating sets
simultaneously without increasing their multiplicity
(\Cref{topology:thicken}), and choose continuous sign coordinates
\(t_i:M\to[0,1]\) with transitions inside these neighborhoods,
equal to one on the inner preimage and zero on the outer exterior
preimage. At most \(n-2\) coordinates are
fractional at any point, so the coordinate map lands in an
\((n-2)\)-dimensional cubical complex. Some coordinate is always one;
equality of coordinate vectors therefore confines the collapse images
of an entire fiber to the same outer ball. Section~\ref{sec:width-map}
triangulates this map and obtains
\[
 C_n=4D+2.
\]

Each label carries a function \(\mu\), which supplies the prescribed weighted mean curvature of
its cut and obeys \(2|\grad\mu|\le a\mu^2+\Pi\), with a coefficient
\(a\in[.60,.63]\) fixed at each stage and a positive
\emph{potential} \(\Pi\); the scalar lower bound required for the cuts
contains a multiple of the sum of these potentials
(\Cref{component:invariant}). A function with constant potential \(p\)
needs length at least \(2\pi/\sqrt{ap}\) to pass from \(+\infty\) to
\(-\infty\). Since the sum of the potentials must be finite, for
infinitely many labels the potentials cannot
have a common positive lower bound, so the corresponding required
separations can be arbitrarily large. We obtain them by stretching the
metric while keeping \(D\) fixed. For a metric \(h\) and a smooth
function \(t\), the weighted scalar curvature is
\begin{equation}\label{intro:weighted-scalar}
 S^{h}_1(t)=\Scal_h+2\Delta_h t-|\dd t|_h^2.
\end{equation}
By \Cref{rounds:torus-curvature}, \(S^{h}_1(t)-|\dd t|_h^2/k\) is the
scalar curvature of \(h+e^{-2t/k}(\dd\theta_1^2+\cdots+\dd\theta_k^2)\)
on \(M\times\T^k\); thus \(S^{h}_1(t)\) is the limit of these
scalar curvatures as \(k\to\infty\).
The local graph deformation of \Cref{step:main} increases \(h\) and
\(t\), retaining a positive quadratic term \(t^2\) in a lower bound for \(S^{h}_1(t)\). Its loss is charged
to the decrease of \(\tr_hh_0\), for a background metric \(h_0\)
fixed throughout a round, plus an arbitrarily small quadratic loss. The
deformation is driven by the forcing function \(F(x,t)\) of \Cref{forcing:derivative-bounds},
whose height derivative \(F_x\) is either at most \(b_0t^2\) or
compensated by its \(t\)-dependence, \(-FF_t\ge P_0F_x\); the constants
of \(F\) are fixed before the local geometry and the error tolerance.
The trace starts at \(n\) and remains positive, so its decrements
telescope independently of the number of overlapping balls.
Rescaling between rounds restores the positive quadratic term while keeping
a fixed positive reserve in original units. Each label stops after
finitely many rounds, and the changes are supported in its original
outer ball, so the construction stabilizes on compact sets
(\Cref{rounds:wide}). Splitting the increments of \(t\) among finitely
many circle warping functions on each compact region converts the resulting bound
into ordinary scalar-curvature control.

The descent uses weighted \(\mu\)-bubbles and torical symmetrization.
To obtain partitions, we retain both sides of each cut as well as its
boundary support, attaching cylinders to the sides to replace their boundaries by complete ends.
The connected components in \Cref{component:category} have product identifications that make later cuts
compatible at intersections (\Cref{component:split}). Besides the usual
boundary stability inequality, this requires a bound on each retained
side. Deleting subsets from a minimizing Caccioppoli set controls the
mean-curvature term at its cylinder interface and gives the required scalar
form bound (\Cref{cut:side-estimate,cut:raw-domination}). The first cut
uses an extra radial label to give compact boundary and side components
before cylinder attachment (see below); after it, a common constant perturbation of the prescribed weighted mean curvature
ensures uniqueness of minimizing Caccioppoli sets on components and their
Euclidean products. Uniqueness then forces the cuts to converge to
the corresponding end base products at infinity
(\Cref{cut:generic-uniqueness,cut:end-convergence}), allowing the
subsequent regularization to make their ends exactly compatible.

The curvature requirement is adapted to the intersections we must
exclude. Each remaining label has a function \(\mu\) and two prescribed
open sets whose closures lie in the interiors of \(\{\mu=+\infty\}\)
and \(\{\mu=-\infty\}\), respectively, where its future sign is fixed;
call the label undecided outside these prescribed open sets.
On a component of base dimension \(m\), we require a positive lower bound
for the stabilized scalar curvature, the scalar curvature of the component
metric warped with all circle factors accumulated so far
(\Cref{rounds:torus-curvature}), only where at least \(m-1\) labels remain
undecided.
A side keeps its dimension and loses the processed label, whereas
a boundary loses both one dimension and that label. Thus every new
required region projects into the old one
(\Cref{component:invariant,component:counts}). This restriction also permits
singular cuts: their singular sets have dimension at most \(m-8\).
Smooth approximations to the remaining functions admit generic
thresholds with at most \(m-8\) equalities at any singular point.
The strict inequalities select signs to force nearby, leaving fewer
than \(m-2\) labels undecided. Enlarging enough of the corresponding
prescribed open sets therefore removes the singular set from the region with at least
\(m-2\) remaining undecided labels, which contains the projections of
all new required regions
(\Cref{local:smooth-approximations,enlargement:global}).

Forcing these additional signs costs potentials concentrated in thin
annuli around singular pieces. On the boundary and on both joined sides,
with matching interface traces, estimates for inverse-square potentials bound each such potential by
a fixed fraction of the reserved gradient form plus a local \(L^2\) error
(\Cref{local:sharp}); the reserved form is a fixed part of the gradient
terms left over in the boundary and side stability forms (\Cref{cut:reserve}). To combine the additions without spending
this reserve once per piece, \Cref{enlargement:maximum} uses compactness:
on a compact region the bounded old potential acts on the limiting
test function, while the shrinking potential acts on the concentrating
remainder. The total form coefficient is therefore bounded by the
larger of the previous and local coefficients, up to an arbitrarily
small error. The extensions in
\Cref{enlargement:global} preserve this bound throughout the component
diagram. The regularization in \Cref{regular:step} can then approximate
the cut closely near the required regions, smooth the interfaces, and use
compatible positive supersolutions to restore the pointwise scalar
bound by adding circle factors. These constructions prove the finite
cutting theorem, \Cref{finite:cut}.

For the countable iteration, the analysis on each component needs one
finite warping function list on the entire complete component: it supplies the
positive injectivity radius and uniform bounds on curvature and all its
covariant derivatives on the whole component and polynomial weighted volume growth
used for the barriers, the uniqueness of minimizing Caccioppoli sets, and the
bounded solution of the elliptic equation (\Cref{component:geometry,countable:locality}). The first, radial cut gives
compact boundary and side components before cylinder attachment; on a
noncompact manifold, alternating signs along a proper radial function
enforce this confinement (\Cref{countable:initial}). Their collapse
images stay in those compact regions, so only finitely many original
outer balls and supports of the logarithms of warping functions meet each image.
\Cref{countable:locality,countable:step} preserve finite lists through
the incident end base diagrams: all other labels remain forced, and their
potentials form a summable list of spatial constants which the
approximations reproduce exactly. The finite cutting argument can
therefore be iterated with summable stage losses. An intersection of
\(n-1\) distinct separating sets would give a complete one-dimensional component.
Its required count is zero, so the positive scalar bound holds
everywhere. The finite warped-product formula then yields a differential
inequality with a negative quadratic upper bound on its arclength universal cover which is impossible on
the whole line (\Cref{countable:dimension-one}). This gives the
intersection bound in \Cref{countable:construction}.

Section~\ref{sec:reconstruction} puts these partitions on one fixed
copy of \(M\). Truncate the analytic cylinders and compress each
inserted finite collar into an arbitrarily thin old collar, preserving
all earlier separating sets and signs. The displacement errors are summable in
the original metric, and the change for a center label is confined to
a fixed-radius original ball about that center. Local finiteness makes
both the collapse maps and the partitions eventually unchanged near
each point. This gives the continuous collapse and locally finite
separating sets used above; locally finite coordinate supports ensure continuity
of the width map in the ordinary polyhedral topology.

The three corollaries are proved in Section~\ref{sec:closed-covers}.
For \Cref{cover:main} in dimensions at least four, apply the theorem
to the normalized lifted metric and realize the coordinate complex
with the uniform simplex metric; local finiteness preserves continuity
and the conversion preserves fibers (\Cref{cover:uniform}). Dimension
two uses Bonnet--Myers. In the noncompact three-dimensional case,
\Cref{cover:level-quotient} makes the bounded level components supplied by
Liokumovich and Wang into entire fibers of a map to the graph obtained by collapsing each connected level component to a point.
The filling-radius corollary follows from Gromov's comparison
\(\operatorname{FillRad}\le\tfrac12\UW_{n-1}\).
For the aspherical obstruction, the width bound on the universal cover
contradicts Gromov's observation that a complete uniformly contractible
\(n\)-manifold, such as the universal cover of a closed aspherical
\(n\)-manifold, admits no continuous map with uniformly bounded fibers
to a polyhedron of dimension less than \(n\).

\paragraph{Organization.}
Section~\ref{sec:width-map} deduces \Cref{thm:main} from the partition
statement, \Cref{topology:partitions}. Sections~\ref{forcing:section},
\ref{step:section} and~\ref{sec:rounds} stretch the metric: the
function with quadratic derivative bounds, the local stretching step, and the rounds of
\Cref{rounds:wide}. Section~\ref{sec:finite-cuts} introduces components,
the analytic invariant and the finite cutting theorem, \Cref{finite:cut}.
Its proof occupies the appendices: component geometry and finite
reconstructions in Appendix~\ref{sec:components}, weighted minimizing cuts
in Appendix~\ref{sec:cut}, exclusion of singular points from the required regions in
Appendices~\ref{sec:local} and~\ref{sec:enlargement-global}, and
regularization in Appendix~\ref{sec:regularization}.
Section~\ref{sec:countable} adapts this construction to countably many
labels and proves the intersection bound, \Cref{countable:construction}.
Section~\ref{sec:reconstruction} proves \Cref{topology:partitions}, and
Section~\ref{sec:closed-covers} proves the three corollaries.

\begin{convention}\label{conv:signs}
The Laplacian is \(\Delta=\Div\grad\), and round spheres have positive
scalar curvature. For a hypersurface with unit normal \(\nu\), its
second fundamental form is
\(B(v,w)=\langle\nabla_v\nu,w\rangle\), and its mean curvature is
\(\tr B\). Integration and divergence use Riemannian densities. Distances in the final fiber bound are always
measured in the original metric \(g\).
\end{convention}
